\documentclass[../main.tex]{subfiles} \begin{document}

\chapter{Smoothed Particle Hydrodynamics}
\label{chap:SPH}

This chapter presents the main numerical method used --- the \emph{Smoothed Particle Method}. In particular,

\begin{compactitem}
	\item section
	\item section
	\item section
\end{compactitem}

\section{Method overview}
\label{sec:SPH_OVERVIEW}

Smoothed Particle Hydrodynamics (SPH) is a general use numerical method first used in astrophysical stellar models \autocite{gingoldSmoothedParticleHydrodynamics1977}. The authors required a solution to a grid resolution problem arising with their use of the finite difference method (FDM).\sidenote{They argued that the treatment of highly distorted continuum would require an extremely fine grid, hence substantial computational cost \autocite[p.~1]{gingoldSmoothedParticleHydrodynamics1977}.} The method has found use in a variety of fields, \viz free-surface flows \autocite{monaghanSimulatingFreeSurface1994a}, solid mechanics \autocite{liberskySmoothParticleHydrodynamics1991}, geomechanics \autocite{buiLagrangianMeshfreeParticles2008}, and is still successfully used in astrophysics \autocite{springelSimulationsFormationEvolution2005}. 

SPH is a particle (Lagrangian) method.\sidenote{In this sense grid/mesh methods such as the \begin{inparaenum} \item finite difference method; \item finite element method; \item finite volume method; \item discontinuous Galerkin method, \end{inparaenum} are called \emph{Eulerian}, since they approximate the evolution in the \emph{current, deformed} configuration. Particle methods like \begin{inparaenum} \item SPH; \item MPS (moving particle semi-implicit method) \autocite{koshizukaMovingParticleSemiImplicitMethod1996}; \item DPD (dissipative particle dynamics) \autocite{hoogerbruggeSimulatingMicroscopicHydrodynamic1992} \end{inparaenum} trace the evolution in the \emph{reference} configuration, and so particle methods are also called \emph{Lagrangian}.}   The continuum is partitioned into arbitrary \emph{material points}, in this context referred to as \emph{particles}, and basic physical properties such as density, velocity, entropy, and energy are assigned to them. Their mutual interaction is then prescribed,\sidenote{The interaction cannot be arbitrary. To capture the physics correctly, the interaction must in fact root in the discretization of the governing equations.} and the state of each particle is advanced in time. 
We stress that \emph{SPH particles} are \emph{not} \textcolor{gray}{elementary} particles in the microscopic sense --- they are typically macroscopic fluid elements. Nonetheless, some intuition from the mechanics of point particles carries over naturally to SPH, as we shall see.

In what follows, we introduce the theoretical foundations of SPH and the specific formulation developed in this work.

\section{Cornerstones of SPH}
\label{sec:CORNERSTONES}

In this section, we outline the standard approach to SPH discretization.

\subsection{Kernel interpolation}
\label{sec:INTERPOLATION}

Let $\Omega \subset \R^d, d = 2, 3$ be the continuous medium and $\Phi: \Rplusnought \times \Omega \to \R$ an arbitrary scalar valued field, \eg density. Suppose further that $\forall t \in \Rplusnought$ the mapping $\vb{x} \mapsto \Phi(t, \vb{x})$ is in $\Loneloc{\Omega}$, then

\begin{equations}[single,!,numberline=all]
\label{eq:DIRAC-PROPERTY}
\Phi\left(t, \vb{x}\right) = \left(\Phi \star \delta\right)\left(t, \vb{x}\right) = \int_{\Omega}\Phi\left(t, \vb{x'}\right)\delta\left(\vb{x} - \vb{x'}\right)\dd{\vb{x'}},
\end{equations}

holds for almost all $\vb{x} \in \Omega$. A common assumption in the SPH literature is to assume $\vb{x} \mapsto \Phi\left(t, \vb{x}\right) \in \Czero{\Omega},$ and so the equality \Cref{eq:DIRAC-PROPERTY} holds pointwise in $\Omega$. To make the presentation simpler, throughout we always assume sufficient regularity of all quantities involved and do not pursue measure-theoretic rigor.

The main point of departure of the SPH method is to replace the Dirac delta distribution in \Cref{eq:DIRAC-PROPERTY} by an \emph{interpolation kernel} $w_h: \Omega \to \R$ \autocite[p.~305]{violeauFluidMechanicsSPH2012}, \ie

\begin{equations}[single,!,numberline=all]
\label{eq:CONT-INTERPOLATION}
\sphc{\Phi}\left(t, \vb{x}\right)\coloneq \int_{\Omega}\Phi\left(t, \vb{x'}\right)w_h\left(\vb{x} - \vb{x'}\right)\dd{\vb{x}'}.
\end{equations}

The map $\Phi \mapsto \sphc{\Phi}$ is called \emph{continuous interpolation}.

- properties of the kernel \autocite[p.~32]{liuSmoothedParticleHydrodynamics2004}
- approximation properties 
- possible choices : Gaussian kernels, $M_n$ splines
- our choice 

\subsection{Particle paradigm}
\label{sec:PARTICLES}

\subsection{Differential operators}
\label{sec:DIFF_OPS}

\section{SPH formulations}
\label{sec:FORMULATIONS}

\subsection{Classical SPH}
\label{sec:WCSPH}

\subsection{Hamiltonian formulation}
\label{sec:HAMILTON}

\subsection{Density-free formulations}
\label{sec:HOPKINS}

\subsubsection{Pressure-entropy}
\label{sec:HOPKINS_PE}
Recall that the entropy of an isentropic flow of an ideal gas has the form\sidenote{This is a test side note.}

\begin{equations}[single,*]
\eta = \eta_0 + c_v \log\left(\frac{P}{\rho^{\gamma}}\right)
\end{equations}
and $\eta = \, \text{const} \,.$ We see that the quantity $A$ defined as \autocite{hopkinsGeneralClassLagrangian2013}

\begin{equations}[single,*]
A = \frac{P}{\rho^{\gamma}},
\end{equations}
is in fact equal to

\begin{equations}[single,*]
A = \exp\left(\frac{\eta - \eta_0}{c_v}\right),
\end{equations}
meaning that the flow is \emph{isentropic,} if and only if the quantity $A$ is \emph{conserved} under the flow. Thus $A$ is an \emph{entropy-like} variable: conserved under isentropic flows.

The choice

\begin{equations}[single,!,numberline=all]
	\label{eq:x_PE}
	x_i = m_i A_i^{\frac{1}{\gamma}},
\end{equations}
gives the pressure $P_i$ as

\begin{equations}[single,!,numberline=all]
	\label{eq:P_PE}
	P_i = \left(\sum_{j=1}^N m_j A_j^{\frac{1}{\gamma}} w_{ij}(h_i)\right)^{\gamma}
\end{equations}
and the evolution equations as

\begin{equations}[single,!,numberline=all]
	\label{eq:v_PE}
	\dv{\vb{v}_i}{t} = - \sum_{j=1}^N m_j \left(A_i A_j\right)^{\frac{1}{\gamma}}\left(\frac{f_{i}P_i}{P_{i}^{\frac{2}{\gamma}}} \nabla_i w_{ij}\left(h_i\right) + \frac{f_{j}P_j}{P^{\frac{2}{\gamma}}_{j}} \nabla_i w_{ij}\left(h_j\right)\right),
\end{equations}
with

\begin{equations}[single,*]
	\label{eq:f_PE}
	f_{i} =  \left(1 + \frac{h_i}{2 n_i} \pdv{n_i}{h_i}\right)^{-1}
\end{equations}

\subsubsection{Pressure-potential temperature}
\label{sec:HOPKINS_theta}


Recall the \emph{potential temperature} $\theta$

\begin{equations}[single,*]
\theta = T \left(\frac{P_0}{P}\right)^{\frac{R}{c_p}},
\end{equations}

where $T$ is the gas's \textcolor{gray}{thermodynamic} temperature, and $P_0$ is the reference pressure. \sidenote{Potential temperature is the thermodynamic temperature the air parcel would possess if \emph{adiabatically} moved from the thermodynamic state $\left(T, P\right)$ to the state $\left(\theta, P_0\right)$.} Using the thermal state equation

\begin{equations}[single,*]
P = R \rho T,
\end{equations}
together with the fact 

\begin{equations}[single,*]
\frac{R}{c_p} = \frac{\gamma - 1}{\gamma},
\end{equations}

we obtain

\begin{equations}[single,*]
\theta = \frac{P}{\rho R} P_0^{\frac{\gamma - 1}{\gamma}} P^{\frac{1 - \gamma}{\gamma}} = \frac{P_0^{\frac{\gamma - 1}{\gamma}}}{R}\left(\frac{P^{\frac{1}{\gamma}}}{\rho}\right) = \frac{P_0^{\frac{\gamma - 1}{\gamma}}}{R}\left(\frac{P}{\rho^{\gamma}}\right)^{\frac{1}{\gamma}} = \frac{P_{0}^{\frac{\gamma - 1}{\gamma}}}{R} A^{\frac{1}{\gamma}},
\end{equations}

so setting 

\begin{equations}[single,*]
\tilde{\theta}_0 \coloneq \frac{P_0^{\frac{\gamma - 1}{\gamma}}}{R}
\end{equations}

we obtain

\begin{equations}[single, !,numberline=all]
	\label{eq:theta_A}
	\theta = \tilde{\theta}_0 A^{\frac{1}{\gamma}}.
\end{equations}

This implies the potential temperature $\theta$ is also \emph{an entropy-like} variable - it is \emph{conserved} under isentropic flows.

Rewriting \ref{eq:x_PE} using \ref{eq:theta_A} we obtain

\begin{equations}[single,!,numberline=all]
\label{eq:x_Ptheta}
x_i = \frac{m_i}{\tilde{\theta}_0} \theta_i.
\end{equations}

Using this quantity the pressure \ref{eq:P_PE} can then be expressed as

\begin{equations}[single,!,numberline=all]
\label{eq:P_Ptheta}
P_i = \left(\sum_{j=1}^N \frac{m_j}{\tilde{\theta_0}} \theta_j w_{ij}\left(h_i\right)\right)^{\gamma}.
\end{equations}

From \ref{eq:v_PE} we obtain

\begin{equations}[single,!,numberline=all]
\label{eq:v_Ptheta}
\dv{\vb{v}_i}{t} = - \sum_{j=1}^N m_j \frac{\theta_i \theta_j}{\tilde{\theta}_0^{2}}\left(\frac{f_{i}P_i}{P_i^{\frac{2}{\gamma}}} \nabla_i w_{ij}\left(h_i\right) + \frac{f_{j}P_j}{P_j^{\frac{2}{\gamma}}} \nabla_i w_{ij}\left(h_j\right)\right),
\end{equations}

with \ref{eq:f_PE} unchanged
\begin{equations}[single,!,numberline=all]
\label{eq:f_Ptheta}
	f_{i} =  \left(1 + \frac{h_i}{2 n_i} \pdv{n_i}{h_i}\right)^{-1}.
\end{equations}

\section{Implementation specifics}
\label{sec:IMPLEMENTATION}

\subsection{Extrapolating boundary conditions}
\label{sec:SPH-BC}

\subsubsection{Inflow \& Outflow}
\label{sec:INFLOW-OUTFLOW}

\subsubsection{Free-slip}
\label{sec:FREE-SLIP}

\subsection{Adaptive smoothing length}
\label{sec:ADAPTIVE_SMOOTHING}
Newton iterations

\subsection{Particle placement}
\label{sec:PLACEMENT}
WitchExpGrid




\end{document}
